\section*{Conclusions}
\label{sec:conclusion}

This paper studies CR-GNN as a controlled biomolecular graph-classification framework rather than as a claim of universal model superiority. The final evidence supports a narrower but stronger conclusion: cross-residual interaction is useful, but residual variants remain the strongest default family in the reported benchmark suite.

The main contribution is therefore not a universally dominant architecture, but a controlled comparison of information-flow mechanisms and residual-routing strategies. Across the biological core, GraphResGNN is strongest on PROTEINS and ENZYMES, whereas NodeResGNN is strongest on DD. Cross-residual variants still provide useful gains over plain stacking on selected datasets, but they do not replace residual baselines as the main point of reference. In this sense, the paper's central finding is comparative and mechanism-oriented: once the training protocol and backbone family are held stable, where residual information is reused matters, and the strongest default remains ordinary residual reuse rather than unrestricted cross exchange.

The routing analysis sharpens that result further. Learnable dense routing remains strongest on the best graph-level residual line, while mild sparse routing is frequently effective on cross-oriented settings. This shows that the practical question is not only whether an additional residual path exists, but how historical information is filtered before reuse. The targeted routing results therefore help explain why the benchmark hierarchy is mixed rather than one-sided.

The scope of this conclusion should remain disciplined. The strongest claims in the paper are supported by the stratified five-fold benchmark and the targeted five-fold ablations, whereas the older sensitivity scans are only supportive. Likewise, the present study addresses controlled graph-classification behavior on medium-scale biomolecular and molecular benchmarks; it does not by itself establish large-scale scalability or direct biological interpretability.

Taken together, these results position CR-GNN as a practical biomolecular graph-classification design space and provide a methodological basis for future work on richer biological inputs, larger benchmarks, and graph-learning settings where residual reuse must be coordinated with stronger domain supervision.
